\section{Introduction}\label{sec:introduction}
The defining recursion for Kazhdan--Lusztig polynomials uses the simple
generators of a Coxeter system, while geometric and categorical
constructions use root labels.  An abstract Bruhat interval poset retains
neither kind of data.

For a Coxeter system $(W,S)$, write $\ell$ for its length function and
$\le$ for Bruhat order. For $u\le b$, the interval
\[
 [u,b]=\{x\in W:u\le x\le b\}
\]
is a finite graded poset of rank $\ell(b)-\ell(u)$, even when $W$ is
infinite. Let $P^W_{u,b}(q)$ denote the usual equal-parameter
Kazhdan--Lusztig polynomial, normalized by $P^W_{u,u}=1$ and
$\deg P^W_{u,b}< (\ell(b)-\ell(u))/2$ for $u<b$.
The combinatorial invariance conjecture asks whether isomorphic interval
posets have the same polynomial. An isomorphism here is only a bijection
preserving and reflecting order; it carries no labels or root data.
The following theorem gives a positive answer in full equal-parameter
generality.

\begin{theorem}[Combinatorial invariance]\label{thm:main}
Let $(W,S)$ and $(W',S')$ be arbitrary Coxeter systems, and let
$u\le b$ in $W$ and $u'\le b'$ in $W'$.
If
\[
 \iota:[u,b]\longrightarrow[u',b']
\]
is an isomorphism of abstract Bruhat interval posets, then their
equal-parameter Kazhdan--Lusztig polynomials satisfy
\[
 P^W_{u,b}(q)=P^{W'}_{u',b'}(q).
\]
\end{theorem}
The statement includes infinite and noncrystallographic groups.
Restricting $\iota$ to subintervals also determines every
polynomial $P_{x,y}$ on the interval; reciprocity then recovers all its
$R$-polynomials (Corollary~\ref{cor:all-subintervals}).

\subsection{History and the missing label data}
Kazhdan and Lusztig introduced these polynomials through the canonical
basis of the Hecke algebra \cite{KL79}. Their geometric interpretation
identifies them with local intersection-cohomology invariants of Schubert
varieties in the Weyl-group setting \cite{KL80}. Thus combinatorial
invariance asks how much of an invariant initially defined by algebra,
and subsequently understood geometrically, is already recorded by a
finite poset. The conjecture is attributed to Lusztig and Dyer
\cite{Delanoy06}. Dyer's published formulation appears in the setting of
finite Coxeter systems in \cite[Section~3.5]{Dyer91}. His reconstruction
of the directed Bruhat graph from the interval order supplies all reflection
edges, including edges that are not covers
\cite[Proposition~3.3]{Dyer91}. Section~\ref{sec:dihedral} gives a closure
proof for arbitrary Coxeter systems.
The reconstruction does not supply root labels.

A principal lower interval is an interval $[1,b]$ starting at the
identity. Special matchings gave an order-theoretic approach to recursion
on these intervals. Du Cloux developed an abstract model for Bruhat
intervals \cite{duCloux00} and proved that isomorphisms of principal lower
intervals preserve all corresponding $P_{x,y}$ when one ambient group
belongs to a class including all finite and affine Coxeter groups
\cite{duCloux03}. Brenti established the corresponding result for
symmetric groups \cite{Brenti04}. Brenti, Caselli, and Marietti proved it
for arbitrary Coxeter systems
\cite[Theorem~7.8 and Corollary~8.4]{BCM06}, as did Delanoy
\cite[Theorem~7.2 and Corollary~7.3]{Delanoy06}.
Their input includes the positions of $x$ and $y$ inside the supplied
principal lower interval; the general conjecture asks for $P_{x,y}$ from the
abstract interval $[x,y]$ alone.

Caselli and Marietti more recently proved that every special matching of
a Bruhat interval in a symmetric group satisfies Brenti's $R$-polynomial
recursion \cite[Theorem~8.6]{CM26}. The same preprint leaves its proposed
algorithm for all type-$A$ intervals conjectural
\cite[Conjecture~8.13 and Proposition~8.15]{CM26}.

For intervals with arbitrary bottom elements, Incitti obtained
short-interval results in classical types \cite{Incitti07}.
Patimo proved invariance of the coefficient of $q$ in finite simply-laced
Weyl groups \cite{Patimo21}. Barkley, Gaetz, and Lam established the same
coefficient invariance for arbitrary Coxeter systems and deduced full
invariance in interval rank at most six
\cite[Corollaries~1.3 and~1.4]{BGL26}.
Combining the flipclass approach with this coefficient result, Esposito,
Marietti, and Stella obtained full invariance through rank ten in finite
Weyl groups and rank twelve in finite type $A$
\cite[Corollary~1.3(ii)]{EMS26}.

Another approach uses hypercube decompositions of Bruhat intervals.
Blundell, Buesing, Davies, Veli\v{c}kovi\'c, and Williamson developed a
recurrence for $P$-polynomials in symmetric groups, guided by analyses of
trained machine-learning models. They proved it for a particular
decomposition defined using the permutation labels of the vertices and
conjectured it for every hypercube decomposition of such an interval
\cite[Theorem~3.7 and Conjecture~3.8]{BBDVW22}.
Barkley and Gaetz proved invariance of the normalized $R$-polynomial
$\widetilde R$ for elementary intervals in symmetric groups
\cite[Theorem~1.6]{BG25Elementary}. They later proved the BBDVW recurrence
for every hypercube decomposition of a principal lower interval $[1,v]$
in a symmetric group \cite[Theorem~1.3]{BG26BBDVW}.

Our approach uses the moment-graph form of the theory.
The fixed-point descriptions of equivariant cohomology developed by
Goresky, Kottwitz, and MacPherson provide the geometric ancestry of its
polynomial congruence conditions \cite{GKM98}.
Braden and MacPherson constructed sheaves that recover intersection
cohomology from moment graphs \cite[Theorems~1.5--1.6]{BM01}.
Soergel's categorification \cite{Soergel07}, Fiebig's moment-graph and
Coxeter-category constructions \cite{FiebigVerma08,Fiebig08,FiebigSurvey},
and the character theorem of Elias and Williamson \cite{EW14} make the
relevant sheaves available for arbitrary Coxeter systems over a suitable
real realization.
Lanini studied invariance through isomorphisms of labeled moment graphs.
Her pullback statement compares Braden--MacPherson sheaves when the graph
isomorphism includes compatible vertexwise automorphisms of a common label
space: they carry edge labels to unit multiples and agree modulo the edge
label at neighboring vertices
\cite[Definitions~2.3--2.4, Lemma~5.1, and Remark~5.1]{Lanini12}.
An abstract Bruhat interval isomorphism does not include these label
automorphisms. We retain the two realizations separately and compare
images in their sheaves through an order on edges.

The reflection-subgroup theory of Deodhar and Dyer, together with Dyer's
reflection-order theory, supplies the combinatorial inputs
\cite{Deodhar89,Dyer87,Dyer90,Dyer91,Dyer93}.
The reflection-path formula already appears in his thesis
\cite[Theorem~6.23]{Dyer87}.
His reconstruction of dihedral coset intersections by four-cycle moves
\cite[Section~3.4]{Dyer91} precedes the relative-interval transport
argument below. Dyer also describes polynomials indexed by initial
sections of a reflection order and interpolating between
$q^{-d/2}P_{u,b}(q)$ and $q^{d/2}P_{u,b}(q^{-1})$, where
$d=\ell(b)-\ell(u)$ \cite[Introduction]{Dyer93}.
His rank-two manuscript contains an integer model of the dihedral graph
and the planar length-sum identity used below, proves scalar congruence
and intersection statements, and sketches localization of projective
modules along cosets
\cite[Sections~1.9--1.11 and~3.16]{DyerRankTwo}.
The proof here combines an exact later-edge count for transported relative
intervals with a bound for the two endpoint images in a sheaf. A
reciprocal comparison forces every such bound to be attained, with
freeness supplied by the same induction.

\subsection{The proof mechanism}
A moment-graph sheaf assigns a graded module to each vertex and edge,
together with restriction maps from an edge's endpoints to its module.
For the sheaf $B(b)$ with top $b$, the vertex module $B^x$ is free over a
real polynomial ring $A$, and its graded rank is $P_{x,b}(q)$.

If $F$ is a set of upward edges at $x$, write
\[
 L_x(F)=\ker\left(B^x\longrightarrow\bigoplus_{e\in F}B^e\right).
\]
These are the vectors whose restrictions vanish along $F$.
Their two extreme cases are known: $L_x(\varnothing)=B^x$, while the
kernel for all upward edges is free with graded rank
$q^{\ell(b)-\ell(x)}P_{x,b}(q^{-1})$.
Intermediate kernels are the objects for which we must establish
freeness.

Choose a reflection order in the second Coxeter system. Use the labels
there to order the edges of the first interval through $\iota$; call the
resulting order a transported schedule. It need not be a reflection order
in the first system. A terminal portion at a vertex consists of all its
edges from a chosen place to the end of this schedule. We prove polynomial
invariance simultaneously with freeness of $L_u(F)$ for every terminal
portion $F$ at the bottom. The induction is on interval rank, so the same
freeness is already available at every strictly higher vertex.

The local step concerns an edge $e:x\to y$. Let $F$ and $G$ be the
strictly later upward edges at $x$ and $y$. Compare their images in the
common edge module:
\[
 I_e=\operatorname{res}_{x,e}L_x(F),\qquad
 J_e=\operatorname{res}_{y,e}L_y(G).
\]
Localization at a plane of root labels decomposes the sheaf into
rank-one dihedral models. Relative dihedral intervals transport through
$\iota$, so their later-edge counts compare these two images.
Assuming only the known freeness of $L_y(G)$, the result is a graded
inclusion
\[
 I_e\subseteq f_eJ_e,\qquad
 \deg f_e=\frac{\ell(y)-\ell(x)-1}{2},
\]
where $\alpha_e$ is the root label of $e$,
$f_e$ is a nonzero homogeneous element of $A/(\alpha_e)$, and
$f_eJ_e$ is free over $A/(\alpha_e)$.
The later global comparison will make this inclusion an equality.

Remove the edge conditions in schedule order. At any real $0<q<1$, the
inclusion bounds the increase in the Hilbert series at $x$ by a positive
multiple of the current kernel's Hilbert series at $y$. The successive
bounds form a product of nonnegative elementary matrices. This product
counts decreasing paths for the reflection order in the second system,
so Dyer's path formula identifies it with that system's normalized
$R$-matrix. Reciprocity in both directions then proves polynomial
equality.

The matrices have nonnegative entries and diagonal entries one. Equality
at the endpoints therefore forces every intermediate bound to be exact,
giving $I_e=f_eJ_e$ for every edge. Adding the bottom conditions back in
reverse order now proves freeness for all terminal kernels and closes the
simultaneous induction.

Section~\ref{sec:background} states the standard inputs and verifies the
realization and grading conventions. Section~\ref{sec:dihedral} proves
graph reconstruction, transport, and the later-edge counts.
Section~\ref{sec:localization} gives
the dihedral models and their local decomposition, and
Section~\ref{sec:edge} obtains the edge-image inclusion.
Section~\ref{sec:induction} carries out reciprocity and the simultaneous
induction.

\subsection{Conventions}
Write $1$ for the identity of a Coxeter group. For comparable elements set
\[
 d(x,y)=\ell(y)-\ell(x).
\]
An upward Bruhat edge is $x\to tx$, where $t$ is a reflection and $tx>x$.
Its label is the positive root of $t$; labels are always on the left.

We use the equal-parameter polynomials with $P_{x,x}=R_{x,x}=1$, and
both polynomials zero when $x\not\le y$.
For $x<y$, $\deg P_{x,y}<d(x,y)/2$.
The $R$-polynomial convention is, for $s\in S$ with $sy<y$,
\[
 R_{x,y}(q)=
 \begin{cases}
 R_{sx,sy}(q),&sx<x,\\
 (q-1)R_{x,sy}(q)+qR_{sx,sy}(q),&sx>x.
 \end{cases}
\]
Together with the degree condition, the reciprocity identity
\begin{equation}\label{eq:reciprocity}
 q^{d(x,y)}P_{x,y}(q^{-1})
 =\sum_{x\le z\le y}R_{x,z}(q)P_{z,y}(q)
\end{equation}
fixes the normalization of $P$ \cite{KL79}.
The same conventions apply to primed systems.
Polynomial rings are over $\mathbb R$, with root labels in degree one.
Ranks of graded free modules record basis degrees with positive
exponents for generators in positive degree.
